\documentclass[11pt,letterpaper]{article}
\usepackage[technical-reference]{cornell-notes}
\usepackage{needspace}

\title{Clause 22: General Utilities Library - Cornell Notes}
\author{}
\date{}
\setDocTitle{Clause 22: General Utilities Library}
\setDocSubtitle{C++ 2024 Cornell Notes}
\setDocOwner{C++ 2024 Cornell Notes}
\setDocAuthor{}
\setDocDate{}
\setCornellCollection{C++ 2024 Cornell Notes}
\setCornellUnitType{Clause}
\setCornellUnitNumber{22}
\setCornellUnitTitle{General Utilities Library}
\setCornellCompanionLabel{C++ Technical Reference}
\hypersetup{pdftitle={Clause 22: General Utilities Library - Cornell Notes},pdfsubject={C++ technical study notes},pdfauthor={}}

\begin{document}
\maketitle
\begin{CornellReferenceOrientationBox}
\textbf{Purpose.} Collects broadly reusable facilities: utilities, tuples, sum types, expected values, bitsets, function wrappers, binders, hashing, type indexing, clocks and execution support, formatting, and other general components.
\par\textbf{Role.} A cross-cutting toolbox shared by applications and other library clauses.
\par\textbf{Principal dependencies.} Uses the library-wide rules of Clause 16 together with the applicable core-language concepts and supporting library clauses.
\par\textbf{Primary audience.} all library users and implementers.
\par\textbf{Major subject groups.} General; Utility components; Pairs; Tuples; Optional objects; Variants; Storage for any type; Expected objects; Bitsets; Function objects; Class type\_\allowbreak{}index; Execution policies; and related subclauses.
\end{CornellReferenceOrientationBox}
\section{Learning Objectives}
\begin{itemize}[label=\textcolor{CornellReferenceTeal}{$\blacktriangleright$}]
\item Define the governing ideas of General and apply its stated contract at a call site.
\item Identify the governing ideas of Utility components and apply its stated contract at a call site.
\item Distinguish the governing ideas of Pairs and apply its stated contract at a call site.
\item Explain the governing ideas of Tuples and apply its stated contract at a call site.
\item Trace the governing ideas of Optional objects and apply its stated contract at a call site.
\item Apply the governing ideas of Variants and apply its stated contract at a call site.
\item Analyze the governing ideas of Storage for any type and apply its stated contract at a call site.
\item Evaluate the governing ideas of Expected objects and apply its stated contract at a call site.
\item Diagnose the governing ideas of Bitsets and apply its stated contract at a call site.
\item Compare the governing ideas of Function objects and apply its stated contract at a call site.
\item Verify the governing ideas of Class type\_\allowbreak{}index and apply its stated contract at a call site.
\item Relate the governing ideas of Execution policies and apply its stated contract at a call site.
\item Classify the governing ideas of Primitive numeric conversions and apply its stated contract at a call site.
\item Justify the governing ideas of Formatting and apply its stated contract at a call site.
\item Audit the governing ideas of Bit manipulation and apply its stated contract at a call site.
\item Summarize the governing ideas of Header <utility> synopsis and apply its stated contract at a call site.
\end{itemize}
\section{Normative Structure Map}
\small\begin{longtable}{@{}>{\RaggedRight\arraybackslash}p{0.13\textwidth}>{\RaggedRight\arraybackslash}p{0.27\textwidth}>{\RaggedRight\arraybackslash}p{0.19\textwidth}>{\RaggedRight\arraybackslash}p{0.31\textwidth}@{}}
\toprule\textbf{Subclause} & \textbf{Subject} & \textbf{Rule category} & \textbf{Key dependencies / entities}\\\midrule\endfirsthead
\toprule\textbf{Subclause} & \textbf{Subject} & \textbf{Rule category} & \textbf{Key dependencies / entities}\\\midrule\endhead
\bottomrule\endfoot
22.1 & General & library semantics & Uses the library-wide rules of Clause 16 together with the applicable core-language concepts and supporting library clauses.\\
22.2 & Utility components & library semantics & Header <utility> synopsis, swap, exchange, Forward/move helpers, and related subclauses\\
22.3 & Pairs & library semantics & General, Class template pair, Specialized algorithms, Tuple-like access to pair, and related subclauses\\
22.4 & Tuples & library semantics & General, Header <tuple> synopsis, Concept tuple-like, Class template tuple, and related subclauses\\
22.5 & Optional objects & library semantics & General, Header <optional> synopsis, Class template optional, No-value state indicator, and related subclauses\\
22.6 & Variants & library semantics & General, Header <variant> synopsis, Class template variant, variant helper classes, and related subclauses\\
22.7 & Storage for any type & library semantics & General, Header <any> synopsis, Class bad\_\allowbreak{}any\_\allowbreak{}cast, Class any, and related subclauses\\
22.8 & Expected objects & library semantics & General, Header <expected> synopsis, Class template unexpected, Class template bad\_\allowbreak{}expected\_\allowbreak{}access, and related subclauses\\
22.9 & Bitsets & library semantics & Header <bitset> synopsis, Class template bitset, bitset hash support, bitset operators\\
22.10 & Function objects & operation semantics & General, Header <functional> synopsis, Definitions, Requirements, and related subclauses\\
22.11 & Class type\_\allowbreak{}index & library semantics & Header <typeindex> synopsis, type\_\allowbreak{}index overview, type\_\allowbreak{}index members, Hash support\\
22.12 & Execution policies & library semantics & General, Header <execution> synopsis, Execution policy type trait, Sequenced execution policy, and related subclauses\\
22.13 & Primitive numeric conversions & library semantics & Header <charconv> synopsis, Primitive numeric output conversion, Primitive numeric input conversion\\
22.14 & Formatting & library semantics & Header <format> synopsis, Format string, Error reporting, Class template basic\_\allowbreak{}format\_\allowbreak{}string, and related subclauses\\
22.15 & Bit manipulation & library semantics & General, Header <bit> synopsis, Function template bit\_\allowbreak{}cast, byteswap, and related subclauses\\
\end{longtable}\normalsize
\begin{CornellReferenceNormativeBox}A heading maps the clause; it does not replace the detailed conditions. For each operation, read requirements, effects, result, exceptions, complexity, remarks, and notes with their stated normative force.\end{CornellReferenceNormativeBox}
\subsection*{Rule Classification Guide}
\small\begin{longtable}{@{}>{\RaggedRight\arraybackslash}p{0.25\textwidth}>{\RaggedRight\arraybackslash}p{0.67\textwidth}@{}}\toprule\textbf{Label or category} & \textbf{How to read it}\\\midrule\endfirsthead\toprule\textbf{Label or category} & \textbf{How to read it (continued)}\\\midrule\endhead\bottomrule\endfoot
Constraints & Conditions controlling whether a declaration, overload, or specialization participates in the relevant context.\\
Mandates & Requirements checked for the described declaration or specialization; failure has the stated well-formedness consequence.\\
Preconditions & Obligations the caller establishes before evaluation; violating one forfeits the operation's contract.\\
Effects / Returns / Postconditions & Separate the state change, produced result, and state guaranteed after successful completion.\\
Throws / Error conditions & State how failure is reported and which exceptional paths are permitted or required.\\
Complexity & Bounds or exact counts for the operations named by the specification.\\
Remarks / Recommended practice & Remarks refine applicability or participation; recommended practice guides implementations without silently becoming a program requirement.\\
Exposition-only & A device for describing semantics, not a public name on which a program can depend.\\
\end{longtable}\normalsize
\section{Cornell Cue-and-Notes}
\Needspace{8\baselineskip}
\subsection*{22.1 General}
\begin{longtable}{@{}>{\RaggedRight\arraybackslash\columncolor{CornellReferenceCueGray}}p{0.28\textwidth}>{\RaggedRight\arraybackslash}p{0.64\textwidth}@{}}
\rowcolor{CornellReferenceNavy}\color{white}\textbf{Cue / recall prompt} & \color{white}\textbf{Notes}\\\endfirsthead
\rowcolor{CornellReferenceNavy}\color{white}\textbf{Cue / recall prompt} & \color{white}\textbf{Notes (continued)}\\\endhead
\arrayrulecolor{CornellReferenceLineGray}
\textbf{What contract governs General?}\\[-1mm]{\scriptsize\CornellClauseRef{22.1}} & Frames the terminology, applicability, and relationships needed to interpret General; detailed rules remain in the following subclauses.\\\hline
\end{longtable}
\begin{CornellReferenceSummaryBox}General supplies the library semantics portion of this clause. The key review move is to connect its local wording to the clause-wide model, then classify each condition by normative force before relying on it.\end{CornellReferenceSummaryBox}
\Needspace{8\baselineskip}
\subsection*{22.2 Utility components}
\begin{longtable}{@{}>{\RaggedRight\arraybackslash\columncolor{CornellReferenceCueGray}}p{0.28\textwidth}>{\RaggedRight\arraybackslash}p{0.64\textwidth}@{}}
\rowcolor{CornellReferenceNavy}\color{white}\textbf{Cue / recall prompt} & \color{white}\textbf{Notes}\\\endfirsthead
\rowcolor{CornellReferenceNavy}\color{white}\textbf{Cue / recall prompt} & \color{white}\textbf{Notes (continued)}\\\endhead
\arrayrulecolor{CornellReferenceLineGray}
\textbf{What contract governs Header <utility> synopsis?}\\[-1mm]{\scriptsize\CornellClauseRef{22.2.1}} & Catalogs the declarations made available by Header <utility> synopsis. The synopsis is an interface map; the detailed specifications supply constraints and behavior.\\\hline
\textbf{What contract governs swap?}\\[-1mm]{\scriptsize\CornellClauseRef{22.2.2}} & Governs state transfer or exchange for swap; check self-operation, validity after movement, exception behavior, and invalidation.\\\hline
\textbf{What contract governs exchange?}\\[-1mm]{\scriptsize\CornellClauseRef{22.2.3}} & Specifies the facility family exchange. Read its declarations together with mandates, preconditions, effects, results, exceptions, complexity, and lifetime rules.\\\hline
\textbf{What contract governs Forward/move helpers?}\\[-1mm]{\scriptsize\CornellClauseRef{22.2.4}} & Specifies the facility family Forward/move helpers. Read its declarations together with mandates, preconditions, effects, results, exceptions, complexity, and lifetime rules.\\\hline
\textbf{What contract governs Function template as\_\allowbreak{}const?}\\[-1mm]{\scriptsize\CornellClauseRef{22.2.5}} & Specifies the facility family Function template as\_\allowbreak{}const. Read its declarations together with mandates, preconditions, effects, results, exceptions, complexity, and lifetime rules.\\\hline
\textbf{What contract governs Function template declval?}\\[-1mm]{\scriptsize\CornellClauseRef{22.2.6}} & Specifies the facility family Function template declval. Read its declarations together with mandates, preconditions, effects, results, exceptions, complexity, and lifetime rules.\\\hline
\textbf{What contract governs Integer comparison functions?}\\[-1mm]{\scriptsize\CornellClauseRef{22.2.7}} & Specifies the facility family Integer comparison functions. Read its declarations together with mandates, preconditions, effects, results, exceptions, complexity, and lifetime rules.\\\hline
\textbf{What contract governs Function template to\_\allowbreak{}underlying?}\\[-1mm]{\scriptsize\CornellClauseRef{22.2.8}} & Specifies the facility family Function template to\_\allowbreak{}underlying. Read its declarations together with mandates, preconditions, effects, results, exceptions, complexity, and lifetime rules.\\\hline
\textbf{What contract governs Function unreachable?}\\[-1mm]{\scriptsize\CornellClauseRef{22.2.9}} & Specifies the facility family Function unreachable. Read its declarations together with mandates, preconditions, effects, results, exceptions, complexity, and lifetime rules.\\\hline
\end{longtable}
\begin{CornellReferenceSummaryBox}Utility components supplies the library semantics portion of this clause. The key review move is to connect its local wording to the clause-wide model, then classify each condition by normative force before relying on it.\end{CornellReferenceSummaryBox}
\Needspace{8\baselineskip}
\subsection*{22.3 Pairs}
\begin{longtable}{@{}>{\RaggedRight\arraybackslash\columncolor{CornellReferenceCueGray}}p{0.28\textwidth}>{\RaggedRight\arraybackslash}p{0.64\textwidth}@{}}
\rowcolor{CornellReferenceNavy}\color{white}\textbf{Cue / recall prompt} & \color{white}\textbf{Notes}\\\endfirsthead
\rowcolor{CornellReferenceNavy}\color{white}\textbf{Cue / recall prompt} & \color{white}\textbf{Notes (continued)}\\\endhead
\arrayrulecolor{CornellReferenceLineGray}
\textbf{What contract governs General?}\\[-1mm]{\scriptsize\CornellClauseRef{22.3.1}} & Frames the terminology, applicability, and relationships needed to interpret General; detailed rules remain in the following subclauses.\\\hline
\textbf{What contract governs Class template pair?}\\[-1mm]{\scriptsize\CornellClauseRef{22.3.2}} & Specifies the facility family Class template pair. Read its declarations together with mandates, preconditions, effects, results, exceptions, complexity, and lifetime rules.\\\hline
\textbf{What does Specialized algorithms guarantee?}\\[-1mm]{\scriptsize\CornellClauseRef{22.3.3}} & Defines the observable result of Specialized algorithms together with applicable iterator-domain, ordering, complexity, invalidation, and exception conditions.\\\hline
\textbf{What contract governs Tuple-like access to pair?}\\[-1mm]{\scriptsize\CornellClauseRef{22.3.4}} & Specifies the facility family Tuple-like access to pair. Read its declarations together with mandates, preconditions, effects, results, exceptions, complexity, and lifetime rules.\\\hline
\textbf{What contract governs Piecewise construction?}\\[-1mm]{\scriptsize\CornellClauseRef{22.3.5}} & Specifies the facility family Piecewise construction. Read its declarations together with mandates, preconditions, effects, results, exceptions, complexity, and lifetime rules.\\\hline
\end{longtable}
\begin{CornellReferenceSummaryBox}Pairs supplies the library semantics portion of this clause. The key review move is to connect its local wording to the clause-wide model, then classify each condition by normative force before relying on it.\end{CornellReferenceSummaryBox}
\Needspace{8\baselineskip}
\subsection*{22.4 Tuples}
\begin{longtable}{@{}>{\RaggedRight\arraybackslash\columncolor{CornellReferenceCueGray}}p{0.28\textwidth}>{\RaggedRight\arraybackslash}p{0.64\textwidth}@{}}
\rowcolor{CornellReferenceNavy}\color{white}\textbf{Cue / recall prompt} & \color{white}\textbf{Notes}\\\endfirsthead
\rowcolor{CornellReferenceNavy}\color{white}\textbf{Cue / recall prompt} & \color{white}\textbf{Notes (continued)}\\\endhead
\arrayrulecolor{CornellReferenceLineGray}
\textbf{What contract governs General?}\\[-1mm]{\scriptsize\CornellClauseRef{22.4.1}} & Frames the terminology, applicability, and relationships needed to interpret General; detailed rules remain in the following subclauses.\\\hline
\textbf{What contract governs Header <tuple> synopsis?}\\[-1mm]{\scriptsize\CornellClauseRef{22.4.2}} & Catalogs the declarations made available by Header <tuple> synopsis. The synopsis is an interface map; the detailed specifications supply constraints and behavior.\\\hline
\textbf{What contract governs Concept tuple-like?}\\[-1mm]{\scriptsize\CornellClauseRef{22.4.3}} & Defines or applies the generic requirement Concept tuple-like; syntactic satisfaction must be paired with every stated semantic and equality-preservation expectation.\\\hline
\textbf{What contract governs Class template tuple?}\\[-1mm]{\scriptsize\CornellClauseRef{22.4.4}} & Specifies the facility family Class template tuple. Read its declarations together with mandates, preconditions, effects, results, exceptions, complexity, and lifetime rules.\\\hline
\textbf{What contract governs Tuple creation functions?}\\[-1mm]{\scriptsize\CornellClauseRef{22.4.5}} & Specifies the facility family Tuple creation functions. Read its declarations together with mandates, preconditions, effects, results, exceptions, complexity, and lifetime rules.\\\hline
\textbf{What contract governs Calling a function with a tuple of arguments?}\\[-1mm]{\scriptsize\CornellClauseRef{22.4.6}} & Specifies the facility family Calling a function with a tuple of arguments. Read its declarations together with mandates, preconditions, effects, results, exceptions, complexity, and lifetime rules.\\\hline
\textbf{What contract governs Tuple helper classes?}\\[-1mm]{\scriptsize\CornellClauseRef{22.4.7}} & Specifies the facility family Tuple helper classes. Read its declarations together with mandates, preconditions, effects, results, exceptions, complexity, and lifetime rules.\\\hline
\textbf{What contract governs Element access?}\\[-1mm]{\scriptsize\CornellClauseRef{22.4.8}} & Specifies the facility family Element access. Read its declarations together with mandates, preconditions, effects, results, exceptions, complexity, and lifetime rules.\\\hline
\textbf{What contract governs Relational operators?}\\[-1mm]{\scriptsize\CornellClauseRef{22.4.9}} & Specifies the facility family Relational operators. Read its declarations together with mandates, preconditions, effects, results, exceptions, complexity, and lifetime rules.\\\hline
\textbf{What contract governs common\_\allowbreak{}reference related specializations?}\\[-1mm]{\scriptsize\CornellClauseRef{22.4.10}} & Specifies the facility family common\_\allowbreak{}reference related specializations. Read its declarations together with mandates, preconditions, effects, results, exceptions, complexity, and lifetime rules.\\\hline
\textbf{What contract governs Tuple traits?}\\[-1mm]{\scriptsize\CornellClauseRef{22.4.11}} & Specifies the facility family Tuple traits. Read its declarations together with mandates, preconditions, effects, results, exceptions, complexity, and lifetime rules.\\\hline
\textbf{What does Tuple specialized algorithms guarantee?}\\[-1mm]{\scriptsize\CornellClauseRef{22.4.12}} & Defines the observable result of Tuple specialized algorithms together with applicable iterator-domain, ordering, complexity, invalidation, and exception conditions.\\\hline
\end{longtable}
\begin{CornellReferenceSummaryBox}Tuples supplies the library semantics portion of this clause. The key review move is to connect its local wording to the clause-wide model, then classify each condition by normative force before relying on it.\end{CornellReferenceSummaryBox}
\Needspace{5\baselineskip}\paragraph{Detailed coverage ledger.}
\begin{description}[style=nextline,leftmargin=2.8cm,font=\normalfont\bfseries\color{CornellReferenceTeal}]
\item[22.4.4.1] \textbf{Construction.} Specifies the facility family Construction. Read its declarations together with mandates, preconditions, effects, results, exceptions, complexity, and lifetime rules.
\item[22.4.4.2] \textbf{Assignment.} Governs state transfer or exchange for Assignment; check self-operation, validity after movement, exception behavior, and invalidation.
\item[22.4.4.3] \textbf{swap.} Governs state transfer or exchange for swap; check self-operation, validity after movement, exception behavior, and invalidation.
\end{description}
\Needspace{8\baselineskip}
\subsection*{22.5 Optional objects}
\begin{longtable}{@{}>{\RaggedRight\arraybackslash\columncolor{CornellReferenceCueGray}}p{0.28\textwidth}>{\RaggedRight\arraybackslash}p{0.64\textwidth}@{}}
\rowcolor{CornellReferenceNavy}\color{white}\textbf{Cue / recall prompt} & \color{white}\textbf{Notes}\\\endfirsthead
\rowcolor{CornellReferenceNavy}\color{white}\textbf{Cue / recall prompt} & \color{white}\textbf{Notes (continued)}\\\endhead
\arrayrulecolor{CornellReferenceLineGray}
\textbf{What contract governs General?}\\[-1mm]{\scriptsize\CornellClauseRef{22.5.1}} & Frames the terminology, applicability, and relationships needed to interpret General; detailed rules remain in the following subclauses.\\\hline
\textbf{What contract governs Header <optional> synopsis?}\\[-1mm]{\scriptsize\CornellClauseRef{22.5.2}} & Catalogs the declarations made available by Header <optional> synopsis. The synopsis is an interface map; the detailed specifications supply constraints and behavior.\\\hline
\textbf{What contract governs Class template optional?}\\[-1mm]{\scriptsize\CornellClauseRef{22.5.3}} & Specifies the facility family Class template optional. Read its declarations together with mandates, preconditions, effects, results, exceptions, complexity, and lifetime rules.\\\hline
\textbf{What contract governs No-value state indicator?}\\[-1mm]{\scriptsize\CornellClauseRef{22.5.4}} & Specifies the facility family No-value state indicator. Read its declarations together with mandates, preconditions, effects, results, exceptions, complexity, and lifetime rules.\\\hline
\textbf{What contract governs Class bad\_\allowbreak{}optional\_\allowbreak{}access?}\\[-1mm]{\scriptsize\CornellClauseRef{22.5.5}} & Specifies the facility family Class bad\_\allowbreak{}optional\_\allowbreak{}access. Read its declarations together with mandates, preconditions, effects, results, exceptions, complexity, and lifetime rules.\\\hline
\textbf{What contract governs Relational operators?}\\[-1mm]{\scriptsize\CornellClauseRef{22.5.6}} & Specifies the facility family Relational operators. Read its declarations together with mandates, preconditions, effects, results, exceptions, complexity, and lifetime rules.\\\hline
\textbf{What contract governs Comparison with nullopt?}\\[-1mm]{\scriptsize\CornellClauseRef{22.5.7}} & Specifies the facility family Comparison with nullopt. Read its declarations together with mandates, preconditions, effects, results, exceptions, complexity, and lifetime rules.\\\hline
\textbf{What contract governs Comparison with T?}\\[-1mm]{\scriptsize\CornellClauseRef{22.5.8}} & Specifies the facility family Comparison with T. Read its declarations together with mandates, preconditions, effects, results, exceptions, complexity, and lifetime rules.\\\hline
\textbf{What does Specialized algorithms guarantee?}\\[-1mm]{\scriptsize\CornellClauseRef{22.5.9}} & Defines the observable result of Specialized algorithms together with applicable iterator-domain, ordering, complexity, invalidation, and exception conditions.\\\hline
\textbf{What contract governs Hash support?}\\[-1mm]{\scriptsize\CornellClauseRef{22.5.10}} & Specifies the facility family Hash support. Read its declarations together with mandates, preconditions, effects, results, exceptions, complexity, and lifetime rules.\\\hline
\end{longtable}
\begin{CornellReferenceSummaryBox}Optional objects supplies the library semantics portion of this clause. The key review move is to connect its local wording to the clause-wide model, then classify each condition by normative force before relying on it.\end{CornellReferenceSummaryBox}
\Needspace{5\baselineskip}\paragraph{Detailed coverage ledger.}
\begin{description}[style=nextline,leftmargin=2.8cm,font=\normalfont\bfseries\color{CornellReferenceTeal}]
\item[22.5.3.1] \textbf{General.} Frames the terminology, applicability, and relationships needed to interpret General; detailed rules remain in the following subclauses.
\item[22.5.3.2] \textbf{Constructors.} Specifies how Constructors establishes object state, including participation constraints, initialization effects, and exceptional exits.
\item[22.5.3.3] \textbf{Destructor.} Specifies how Destructor ends the object's managed state and which cleanup or termination consequences apply.
\item[22.5.3.4] \textbf{Assignment.} Governs state transfer or exchange for Assignment; check self-operation, validity after movement, exception behavior, and invalidation.
\item[22.5.3.5] \textbf{Swap.} Governs state transfer or exchange for Swap; check self-operation, validity after movement, exception behavior, and invalidation.
\item[22.5.3.6] \textbf{Observers.} Specifies the facility family Observers. Read its declarations together with mandates, preconditions, effects, results, exceptions, complexity, and lifetime rules.
\item[22.5.3.7] \textbf{Monadic operations.} Defines the observable result of Monadic operations together with applicable iterator-domain, ordering, complexity, invalidation, and exception conditions.
\item[22.5.3.8] \textbf{Modifiers.} Specifies the facility family Modifiers. Read its declarations together with mandates, preconditions, effects, results, exceptions, complexity, and lifetime rules.
\end{description}
\Needspace{8\baselineskip}
\subsection*{22.6 Variants}
\begin{longtable}{@{}>{\RaggedRight\arraybackslash\columncolor{CornellReferenceCueGray}}p{0.28\textwidth}>{\RaggedRight\arraybackslash}p{0.64\textwidth}@{}}
\rowcolor{CornellReferenceNavy}\color{white}\textbf{Cue / recall prompt} & \color{white}\textbf{Notes}\\\endfirsthead
\rowcolor{CornellReferenceNavy}\color{white}\textbf{Cue / recall prompt} & \color{white}\textbf{Notes (continued)}\\\endhead
\arrayrulecolor{CornellReferenceLineGray}
\textbf{What contract governs General?}\\[-1mm]{\scriptsize\CornellClauseRef{22.6.1}} & Frames the terminology, applicability, and relationships needed to interpret General; detailed rules remain in the following subclauses.\\\hline
\textbf{What contract governs Header <variant> synopsis?}\\[-1mm]{\scriptsize\CornellClauseRef{22.6.2}} & Catalogs the declarations made available by Header <variant> synopsis. The synopsis is an interface map; the detailed specifications supply constraints and behavior.\\\hline
\textbf{What contract governs Class template variant?}\\[-1mm]{\scriptsize\CornellClauseRef{22.6.3}} & Specifies the facility family Class template variant. Read its declarations together with mandates, preconditions, effects, results, exceptions, complexity, and lifetime rules.\\\hline
\textbf{What contract governs variant helper classes?}\\[-1mm]{\scriptsize\CornellClauseRef{22.6.4}} & Specifies the facility family variant helper classes. Read its declarations together with mandates, preconditions, effects, results, exceptions, complexity, and lifetime rules.\\\hline
\textbf{What contract governs Value access?}\\[-1mm]{\scriptsize\CornellClauseRef{22.6.5}} & Specifies the facility family Value access. Read its declarations together with mandates, preconditions, effects, results, exceptions, complexity, and lifetime rules.\\\hline
\textbf{What contract governs Relational operators?}\\[-1mm]{\scriptsize\CornellClauseRef{22.6.6}} & Specifies the facility family Relational operators. Read its declarations together with mandates, preconditions, effects, results, exceptions, complexity, and lifetime rules.\\\hline
\textbf{What contract governs Visitation?}\\[-1mm]{\scriptsize\CornellClauseRef{22.6.7}} & Specifies the facility family Visitation. Read its declarations together with mandates, preconditions, effects, results, exceptions, complexity, and lifetime rules.\\\hline
\textbf{What contract governs Class monostate?}\\[-1mm]{\scriptsize\CornellClauseRef{22.6.8}} & Specifies the facility family Class monostate. Read its declarations together with mandates, preconditions, effects, results, exceptions, complexity, and lifetime rules.\\\hline
\textbf{What contract governs monostate relational operators?}\\[-1mm]{\scriptsize\CornellClauseRef{22.6.9}} & Specifies the facility family monostate relational operators. Read its declarations together with mandates, preconditions, effects, results, exceptions, complexity, and lifetime rules.\\\hline
\textbf{What does Specialized algorithms guarantee?}\\[-1mm]{\scriptsize\CornellClauseRef{22.6.10}} & Defines the observable result of Specialized algorithms together with applicable iterator-domain, ordering, complexity, invalidation, and exception conditions.\\\hline
\textbf{What contract governs Class bad\_\allowbreak{}variant\_\allowbreak{}access?}\\[-1mm]{\scriptsize\CornellClauseRef{22.6.11}} & Specifies the facility family Class bad\_\allowbreak{}variant\_\allowbreak{}access. Read its declarations together with mandates, preconditions, effects, results, exceptions, complexity, and lifetime rules.\\\hline
\textbf{What contract governs Hash support?}\\[-1mm]{\scriptsize\CornellClauseRef{22.6.12}} & Specifies the facility family Hash support. Read its declarations together with mandates, preconditions, effects, results, exceptions, complexity, and lifetime rules.\\\hline
\end{longtable}
\begin{CornellReferenceSummaryBox}Variants supplies the library semantics portion of this clause. The key review move is to connect its local wording to the clause-wide model, then classify each condition by normative force before relying on it.\end{CornellReferenceSummaryBox}
\Needspace{5\baselineskip}\paragraph{Detailed coverage ledger.}
\begin{description}[style=nextline,leftmargin=2.8cm,font=\normalfont\bfseries\color{CornellReferenceTeal}]
\item[22.6.3.1] \textbf{General.} Frames the terminology, applicability, and relationships needed to interpret General; detailed rules remain in the following subclauses.
\item[22.6.3.2] \textbf{Constructors.} Specifies how Constructors establishes object state, including participation constraints, initialization effects, and exceptional exits.
\item[22.6.3.3] \textbf{Destructor.} Specifies how Destructor ends the object's managed state and which cleanup or termination consequences apply.
\item[22.6.3.4] \textbf{Assignment.} Governs state transfer or exchange for Assignment; check self-operation, validity after movement, exception behavior, and invalidation.
\item[22.6.3.5] \textbf{Modifiers.} Specifies the facility family Modifiers. Read its declarations together with mandates, preconditions, effects, results, exceptions, complexity, and lifetime rules.
\item[22.6.3.6] \textbf{Value status.} Specifies the facility family Value status. Read its declarations together with mandates, preconditions, effects, results, exceptions, complexity, and lifetime rules.
\item[22.6.3.7] \textbf{Swap.} Governs state transfer or exchange for Swap; check self-operation, validity after movement, exception behavior, and invalidation.
\end{description}
\Needspace{8\baselineskip}
\subsection*{22.7 Storage for any type}
\begin{longtable}{@{}>{\RaggedRight\arraybackslash\columncolor{CornellReferenceCueGray}}p{0.28\textwidth}>{\RaggedRight\arraybackslash}p{0.64\textwidth}@{}}
\rowcolor{CornellReferenceNavy}\color{white}\textbf{Cue / recall prompt} & \color{white}\textbf{Notes}\\\endfirsthead
\rowcolor{CornellReferenceNavy}\color{white}\textbf{Cue / recall prompt} & \color{white}\textbf{Notes (continued)}\\\endhead
\arrayrulecolor{CornellReferenceLineGray}
\textbf{What contract governs General?}\\[-1mm]{\scriptsize\CornellClauseRef{22.7.1}} & Frames the terminology, applicability, and relationships needed to interpret General; detailed rules remain in the following subclauses.\\\hline
\textbf{What contract governs Header <any> synopsis?}\\[-1mm]{\scriptsize\CornellClauseRef{22.7.2}} & Catalogs the declarations made available by Header <any> synopsis. The synopsis is an interface map; the detailed specifications supply constraints and behavior.\\\hline
\textbf{What contract governs Class bad\_\allowbreak{}any\_\allowbreak{}cast?}\\[-1mm]{\scriptsize\CornellClauseRef{22.7.3}} & Specifies the facility family Class bad\_\allowbreak{}any\_\allowbreak{}cast. Read its declarations together with mandates, preconditions, effects, results, exceptions, complexity, and lifetime rules.\\\hline
\textbf{What contract governs Class any?}\\[-1mm]{\scriptsize\CornellClauseRef{22.7.4}} & Specifies the facility family Class any. Read its declarations together with mandates, preconditions, effects, results, exceptions, complexity, and lifetime rules.\\\hline
\textbf{What contract governs Non-member functions?}\\[-1mm]{\scriptsize\CornellClauseRef{22.7.5}} & Specifies the facility family Non-member functions. Read its declarations together with mandates, preconditions, effects, results, exceptions, complexity, and lifetime rules.\\\hline
\end{longtable}
\begin{CornellReferenceSummaryBox}Storage for any type supplies the library semantics portion of this clause. The key review move is to connect its local wording to the clause-wide model, then classify each condition by normative force before relying on it.\end{CornellReferenceSummaryBox}
\Needspace{5\baselineskip}\paragraph{Detailed coverage ledger.}
\begin{description}[style=nextline,leftmargin=2.8cm,font=\normalfont\bfseries\color{CornellReferenceTeal}]
\item[22.7.4.1] \textbf{General.} Frames the terminology, applicability, and relationships needed to interpret General; detailed rules remain in the following subclauses.
\item[22.7.4.2] \textbf{Construction and destruction.} Specifies the facility family Construction and destruction. Read its declarations together with mandates, preconditions, effects, results, exceptions, complexity, and lifetime rules.
\item[22.7.4.3] \textbf{Assignment.} Governs state transfer or exchange for Assignment; check self-operation, validity after movement, exception behavior, and invalidation.
\item[22.7.4.4] \textbf{Modifiers.} Specifies the facility family Modifiers. Read its declarations together with mandates, preconditions, effects, results, exceptions, complexity, and lifetime rules.
\item[22.7.4.5] \textbf{Observers.} Specifies the facility family Observers. Read its declarations together with mandates, preconditions, effects, results, exceptions, complexity, and lifetime rules.
\end{description}
\Needspace{8\baselineskip}
\subsection*{22.8 Expected objects}
\begin{longtable}{@{}>{\RaggedRight\arraybackslash\columncolor{CornellReferenceCueGray}}p{0.28\textwidth}>{\RaggedRight\arraybackslash}p{0.64\textwidth}@{}}
\rowcolor{CornellReferenceNavy}\color{white}\textbf{Cue / recall prompt} & \color{white}\textbf{Notes}\\\endfirsthead
\rowcolor{CornellReferenceNavy}\color{white}\textbf{Cue / recall prompt} & \color{white}\textbf{Notes (continued)}\\\endhead
\arrayrulecolor{CornellReferenceLineGray}
\textbf{What contract governs General?}\\[-1mm]{\scriptsize\CornellClauseRef{22.8.1}} & Frames the terminology, applicability, and relationships needed to interpret General; detailed rules remain in the following subclauses.\\\hline
\textbf{What contract governs Header <expected> synopsis?}\\[-1mm]{\scriptsize\CornellClauseRef{22.8.2}} & Catalogs the declarations made available by Header <expected> synopsis. The synopsis is an interface map; the detailed specifications supply constraints and behavior.\\\hline
\textbf{What contract governs Class template unexpected?}\\[-1mm]{\scriptsize\CornellClauseRef{22.8.3}} & Specifies the facility family Class template unexpected. Read its declarations together with mandates, preconditions, effects, results, exceptions, complexity, and lifetime rules.\\\hline
\textbf{What contract governs Class template bad\_\allowbreak{}expected\_\allowbreak{}access?}\\[-1mm]{\scriptsize\CornellClauseRef{22.8.4}} & Specifies the facility family Class template bad\_\allowbreak{}expected\_\allowbreak{}access. Read its declarations together with mandates, preconditions, effects, results, exceptions, complexity, and lifetime rules.\\\hline
\textbf{What contract governs Class template specialization bad\_\allowbreak{}expected\_\allowbreak{}access<void>?}\\[-1mm]{\scriptsize\CornellClauseRef{22.8.5}} & Specifies the facility family Class template specialization bad\_\allowbreak{}expected\_\allowbreak{}access<void>. Read its declarations together with mandates, preconditions, effects, results, exceptions, complexity, and lifetime rules.\\\hline
\textbf{What contract governs Class template expected?}\\[-1mm]{\scriptsize\CornellClauseRef{22.8.6}} & Specifies the facility family Class template expected. Read its declarations together with mandates, preconditions, effects, results, exceptions, complexity, and lifetime rules.\\\hline
\textbf{What contract governs Partial specialization of expected for void types?}\\[-1mm]{\scriptsize\CornellClauseRef{22.8.7}} & Specifies the facility family Partial specialization of expected for void types. Read its declarations together with mandates, preconditions, effects, results, exceptions, complexity, and lifetime rules.\\\hline
\end{longtable}
\begin{CornellReferenceSummaryBox}Expected objects supplies the library semantics portion of this clause. The key review move is to connect its local wording to the clause-wide model, then classify each condition by normative force before relying on it.\end{CornellReferenceSummaryBox}
\Needspace{5\baselineskip}\paragraph{Detailed coverage ledger.}
\begin{description}[style=nextline,leftmargin=2.8cm,font=\normalfont\bfseries\color{CornellReferenceTeal}]
\item[22.8.3.1] \textbf{General.} Frames the terminology, applicability, and relationships needed to interpret General; detailed rules remain in the following subclauses.
\item[22.8.3.2] \textbf{Constructors.} Specifies how Constructors establishes object state, including participation constraints, initialization effects, and exceptional exits.
\item[22.8.3.3] \textbf{Observers.} Specifies the facility family Observers. Read its declarations together with mandates, preconditions, effects, results, exceptions, complexity, and lifetime rules.
\item[22.8.3.4] \textbf{Swap.} Governs state transfer or exchange for Swap; check self-operation, validity after movement, exception behavior, and invalidation.
\item[22.8.3.5] \textbf{Equality operator.} Specifies the facility family Equality operator. Read its declarations together with mandates, preconditions, effects, results, exceptions, complexity, and lifetime rules.
\item[22.8.6.1] \textbf{General.} Frames the terminology, applicability, and relationships needed to interpret General; detailed rules remain in the following subclauses.
\item[22.8.6.2] \textbf{Constructors.} Specifies how Constructors establishes object state, including participation constraints, initialization effects, and exceptional exits.
\item[22.8.6.3] \textbf{Destructor.} Specifies how Destructor ends the object's managed state and which cleanup or termination consequences apply.
\item[22.8.6.4] \textbf{Assignment.} Governs state transfer or exchange for Assignment; check self-operation, validity after movement, exception behavior, and invalidation.
\item[22.8.6.5] \textbf{Swap.} Governs state transfer or exchange for Swap; check self-operation, validity after movement, exception behavior, and invalidation.
\item[22.8.6.6] \textbf{Observers.} Specifies the facility family Observers. Read its declarations together with mandates, preconditions, effects, results, exceptions, complexity, and lifetime rules.
\item[22.8.6.7] \textbf{Monadic operations.} Defines the observable result of Monadic operations together with applicable iterator-domain, ordering, complexity, invalidation, and exception conditions.
\item[22.8.6.8] \textbf{Equality operators.} Specifies the facility family Equality operators. Read its declarations together with mandates, preconditions, effects, results, exceptions, complexity, and lifetime rules.
\item[22.8.7.1] \textbf{General.} Frames the terminology, applicability, and relationships needed to interpret General; detailed rules remain in the following subclauses.
\item[22.8.7.2] \textbf{Constructors.} Specifies how Constructors establishes object state, including participation constraints, initialization effects, and exceptional exits.
\item[22.8.7.3] \textbf{Destructor.} Specifies how Destructor ends the object's managed state and which cleanup or termination consequences apply.
\item[22.8.7.4] \textbf{Assignment.} Governs state transfer or exchange for Assignment; check self-operation, validity after movement, exception behavior, and invalidation.
\item[22.8.7.5] \textbf{Swap.} Governs state transfer or exchange for Swap; check self-operation, validity after movement, exception behavior, and invalidation.
\item[22.8.7.6] \textbf{Observers.} Specifies the facility family Observers. Read its declarations together with mandates, preconditions, effects, results, exceptions, complexity, and lifetime rules.
\item[22.8.7.7] \textbf{Monadic operations.} Defines the observable result of Monadic operations together with applicable iterator-domain, ordering, complexity, invalidation, and exception conditions.
\item[22.8.7.8] \textbf{Equality operators.} Specifies the facility family Equality operators. Read its declarations together with mandates, preconditions, effects, results, exceptions, complexity, and lifetime rules.
\end{description}
\Needspace{8\baselineskip}
\subsection*{22.9 Bitsets}
\begin{longtable}{@{}>{\RaggedRight\arraybackslash\columncolor{CornellReferenceCueGray}}p{0.28\textwidth}>{\RaggedRight\arraybackslash}p{0.64\textwidth}@{}}
\rowcolor{CornellReferenceNavy}\color{white}\textbf{Cue / recall prompt} & \color{white}\textbf{Notes}\\\endfirsthead
\rowcolor{CornellReferenceNavy}\color{white}\textbf{Cue / recall prompt} & \color{white}\textbf{Notes (continued)}\\\endhead
\arrayrulecolor{CornellReferenceLineGray}
\textbf{What contract governs Header <bitset> synopsis?}\\[-1mm]{\scriptsize\CornellClauseRef{22.9.1}} & Catalogs the declarations made available by Header <bitset> synopsis. The synopsis is an interface map; the detailed specifications supply constraints and behavior.\\\hline
\textbf{What contract governs Class template bitset?}\\[-1mm]{\scriptsize\CornellClauseRef{22.9.2}} & Specifies the facility family Class template bitset. Read its declarations together with mandates, preconditions, effects, results, exceptions, complexity, and lifetime rules.\\\hline
\textbf{What contract governs bitset hash support?}\\[-1mm]{\scriptsize\CornellClauseRef{22.9.3}} & Specifies the facility family bitset hash support. Read its declarations together with mandates, preconditions, effects, results, exceptions, complexity, and lifetime rules.\\\hline
\textbf{What contract governs bitset operators?}\\[-1mm]{\scriptsize\CornellClauseRef{22.9.4}} & Specifies the facility family bitset operators. Read its declarations together with mandates, preconditions, effects, results, exceptions, complexity, and lifetime rules.\\\hline
\end{longtable}
\begin{CornellReferenceSummaryBox}Bitsets supplies the library semantics portion of this clause. The key review move is to connect its local wording to the clause-wide model, then classify each condition by normative force before relying on it.\end{CornellReferenceSummaryBox}
\Needspace{5\baselineskip}\paragraph{Detailed coverage ledger.}
\begin{description}[style=nextline,leftmargin=2.8cm,font=\normalfont\bfseries\color{CornellReferenceTeal}]
\item[22.9.2.1] \textbf{General.} Frames the terminology, applicability, and relationships needed to interpret General; detailed rules remain in the following subclauses.
\item[22.9.2.2] \textbf{Constructors.} Specifies how Constructors establishes object state, including participation constraints, initialization effects, and exceptional exits.
\item[22.9.2.3] \textbf{Members.} Specifies the facility family Members. Read its declarations together with mandates, preconditions, effects, results, exceptions, complexity, and lifetime rules.
\end{description}
\Needspace{8\baselineskip}
\subsection*{22.10 Function objects}
\begin{longtable}{@{}>{\RaggedRight\arraybackslash\columncolor{CornellReferenceCueGray}}p{0.28\textwidth}>{\RaggedRight\arraybackslash}p{0.64\textwidth}@{}}
\rowcolor{CornellReferenceNavy}\color{white}\textbf{Cue / recall prompt} & \color{white}\textbf{Notes}\\\endfirsthead
\rowcolor{CornellReferenceNavy}\color{white}\textbf{Cue / recall prompt} & \color{white}\textbf{Notes (continued)}\\\endhead
\arrayrulecolor{CornellReferenceLineGray}
\textbf{What contract governs General?}\\[-1mm]{\scriptsize\CornellClauseRef{22.10.1}} & Frames the terminology, applicability, and relationships needed to interpret General; detailed rules remain in the following subclauses.\\\hline
\textbf{What contract governs Header <functional> synopsis?}\\[-1mm]{\scriptsize\CornellClauseRef{22.10.2}} & Catalogs the declarations made available by Header <functional> synopsis. The synopsis is an interface map; the detailed specifications supply constraints and behavior.\\\hline
\textbf{What contract governs Definitions?}\\[-1mm]{\scriptsize\CornellClauseRef{22.10.3}} & Specifies the facility family Definitions. Read its declarations together with mandates, preconditions, effects, results, exceptions, complexity, and lifetime rules.\\\hline
\textbf{Which obligations govern Requirements?}\\[-1mm]{\scriptsize\CornellClauseRef{22.10.4}} & States the admissibility and semantic obligations for Requirements. Separate compile-time participation rules from caller preconditions and runtime effects.\\\hline
\textbf{What contract governs invoke functions?}\\[-1mm]{\scriptsize\CornellClauseRef{22.10.5}} & Specifies the facility family invoke functions. Read its declarations together with mandates, preconditions, effects, results, exceptions, complexity, and lifetime rules.\\\hline
\textbf{What contract governs Class template reference\_\allowbreak{}wrapper?}\\[-1mm]{\scriptsize\CornellClauseRef{22.10.6}} & Specifies the facility family Class template reference\_\allowbreak{}wrapper. Read its declarations together with mandates, preconditions, effects, results, exceptions, complexity, and lifetime rules.\\\hline
\textbf{What contract governs Arithmetic operations?}\\[-1mm]{\scriptsize\CornellClauseRef{22.10.7}} & Defines the observable result of Arithmetic operations together with applicable iterator-domain, ordering, complexity, invalidation, and exception conditions.\\\hline
\textbf{What contract governs Comparisons?}\\[-1mm]{\scriptsize\CornellClauseRef{22.10.8}} & Specifies the facility family Comparisons. Read its declarations together with mandates, preconditions, effects, results, exceptions, complexity, and lifetime rules.\\\hline
\textbf{What contract governs Concept-constrained comparisons?}\\[-1mm]{\scriptsize\CornellClauseRef{22.10.9}} & Defines or applies the generic requirement Concept-constrained comparisons; syntactic satisfaction must be paired with every stated semantic and equality-preservation expectation.\\\hline
\textbf{What contract governs Logical operations?}\\[-1mm]{\scriptsize\CornellClauseRef{22.10.10}} & Defines the observable result of Logical operations together with applicable iterator-domain, ordering, complexity, invalidation, and exception conditions.\\\hline
\textbf{What contract governs Bitwise operations?}\\[-1mm]{\scriptsize\CornellClauseRef{22.10.11}} & Defines the observable result of Bitwise operations together with applicable iterator-domain, ordering, complexity, invalidation, and exception conditions.\\\hline
\textbf{What contract governs Class identity?}\\[-1mm]{\scriptsize\CornellClauseRef{22.10.12}} & Specifies the facility family Class identity. Read its declarations together with mandates, preconditions, effects, results, exceptions, complexity, and lifetime rules.\\\hline
\textbf{What contract governs Function template not\_\allowbreak{}fn?}\\[-1mm]{\scriptsize\CornellClauseRef{22.10.13}} & Specifies the facility family Function template not\_\allowbreak{}fn. Read its declarations together with mandates, preconditions, effects, results, exceptions, complexity, and lifetime rules.\\\hline
\textbf{What contract governs Function templates bind\_\allowbreak{}front and bind\_\allowbreak{}back?}\\[-1mm]{\scriptsize\CornellClauseRef{22.10.14}} & Specifies the facility family Function templates bind\_\allowbreak{}front and bind\_\allowbreak{}back. Read its declarations together with mandates, preconditions, effects, results, exceptions, complexity, and lifetime rules.\\\hline
\textbf{What contract governs Function object binders?}\\[-1mm]{\scriptsize\CornellClauseRef{22.10.15}} & Specifies the facility family Function object binders. Read its declarations together with mandates, preconditions, effects, results, exceptions, complexity, and lifetime rules.\\\hline
\textbf{What contract governs Function template mem\_\allowbreak{}fn?}\\[-1mm]{\scriptsize\CornellClauseRef{22.10.16}} & Specifies the facility family Function template mem\_\allowbreak{}fn. Read its declarations together with mandates, preconditions, effects, results, exceptions, complexity, and lifetime rules.\\\hline
\textbf{What contract governs Polymorphic function wrappers?}\\[-1mm]{\scriptsize\CornellClauseRef{22.10.17}} & Specifies the facility family Polymorphic function wrappers. Read its declarations together with mandates, preconditions, effects, results, exceptions, complexity, and lifetime rules.\\\hline
\textbf{What contract governs Searchers?}\\[-1mm]{\scriptsize\CornellClauseRef{22.10.18}} & Specifies the facility family Searchers. Read its declarations together with mandates, preconditions, effects, results, exceptions, complexity, and lifetime rules.\\\hline
\textbf{What contract governs Class template hash?}\\[-1mm]{\scriptsize\CornellClauseRef{22.10.19}} & Specifies the facility family Class template hash. Read its declarations together with mandates, preconditions, effects, results, exceptions, complexity, and lifetime rules.\\\hline
\end{longtable}
\begin{CornellReferenceSummaryBox}Function objects supplies the operation semantics portion of this clause. The key review move is to connect its local wording to the clause-wide model, then classify each condition by normative force before relying on it.\end{CornellReferenceSummaryBox}
\Needspace{5\baselineskip}\paragraph{Detailed coverage ledger.}
\begin{description}[style=nextline,leftmargin=2.8cm,font=\normalfont\bfseries\color{CornellReferenceTeal}]
\item[22.10.6.1] \textbf{General.} Frames the terminology, applicability, and relationships needed to interpret General; detailed rules remain in the following subclauses.
\item[22.10.6.2] \textbf{Constructors.} Specifies how Constructors establishes object state, including participation constraints, initialization effects, and exceptional exits.
\item[22.10.6.3] \textbf{Assignment.} Governs state transfer or exchange for Assignment; check self-operation, validity after movement, exception behavior, and invalidation.
\item[22.10.6.4] \textbf{Access.} Specifies the facility family Access. Read its declarations together with mandates, preconditions, effects, results, exceptions, complexity, and lifetime rules.
\item[22.10.6.5] \textbf{Invocation.} Specifies the facility family Invocation. Read its declarations together with mandates, preconditions, effects, results, exceptions, complexity, and lifetime rules.
\item[22.10.6.6] \textbf{Helper functions.} Specifies the facility family Helper functions. Read its declarations together with mandates, preconditions, effects, results, exceptions, complexity, and lifetime rules.
\item[22.10.6.7] \textbf{common\_\allowbreak{}reference related specializations.} Specifies the facility family common\_\allowbreak{}reference related specializations. Read its declarations together with mandates, preconditions, effects, results, exceptions, complexity, and lifetime rules.
\item[22.10.7.1] \textbf{General.} Frames the terminology, applicability, and relationships needed to interpret General; detailed rules remain in the following subclauses.
\item[22.10.7.2] \textbf{Class template plus.} Specifies the facility family Class template plus. Read its declarations together with mandates, preconditions, effects, results, exceptions, complexity, and lifetime rules.
\item[22.10.7.3] \textbf{Class template minus.} Specifies the facility family Class template minus. Read its declarations together with mandates, preconditions, effects, results, exceptions, complexity, and lifetime rules.
\item[22.10.7.4] \textbf{Class template multiplies.} Specifies the facility family Class template multiplies. Read its declarations together with mandates, preconditions, effects, results, exceptions, complexity, and lifetime rules.
\item[22.10.7.5] \textbf{Class template divides.} Specifies the facility family Class template divides. Read its declarations together with mandates, preconditions, effects, results, exceptions, complexity, and lifetime rules.
\item[22.10.7.6] \textbf{Class template modulus.} Specifies the facility family Class template modulus. Read its declarations together with mandates, preconditions, effects, results, exceptions, complexity, and lifetime rules.
\item[22.10.7.7] \textbf{Class template negate.} Specifies the facility family Class template negate. Read its declarations together with mandates, preconditions, effects, results, exceptions, complexity, and lifetime rules.
\item[22.10.8.1] \textbf{General.} Frames the terminology, applicability, and relationships needed to interpret General; detailed rules remain in the following subclauses.
\item[22.10.8.2] \textbf{Class template equal\_\allowbreak{}to.} Specifies the facility family Class template equal\_\allowbreak{}to. Read its declarations together with mandates, preconditions, effects, results, exceptions, complexity, and lifetime rules.
\item[22.10.8.3] \textbf{Class template not\_\allowbreak{}equal\_\allowbreak{}to.} Specifies the facility family Class template not\_\allowbreak{}equal\_\allowbreak{}to. Read its declarations together with mandates, preconditions, effects, results, exceptions, complexity, and lifetime rules.
\item[22.10.8.4] \textbf{Class template greater.} Specifies the facility family Class template greater. Read its declarations together with mandates, preconditions, effects, results, exceptions, complexity, and lifetime rules.
\item[22.10.8.5] \textbf{Class template less.} Specifies the facility family Class template less. Read its declarations together with mandates, preconditions, effects, results, exceptions, complexity, and lifetime rules.
\item[22.10.8.6] \textbf{Class template greater\_\allowbreak{}equal.} Specifies the facility family Class template greater\_\allowbreak{}equal. Read its declarations together with mandates, preconditions, effects, results, exceptions, complexity, and lifetime rules.
\item[22.10.8.7] \textbf{Class template less\_\allowbreak{}equal.} Specifies the facility family Class template less\_\allowbreak{}equal. Read its declarations together with mandates, preconditions, effects, results, exceptions, complexity, and lifetime rules.
\item[22.10.8.8] \textbf{Class compare\_\allowbreak{}three\_\allowbreak{}way.} Specifies the facility family Class compare\_\allowbreak{}three\_\allowbreak{}way. Read its declarations together with mandates, preconditions, effects, results, exceptions, complexity, and lifetime rules.
\item[22.10.10.1] \textbf{General.} Frames the terminology, applicability, and relationships needed to interpret General; detailed rules remain in the following subclauses.
\item[22.10.10.2] \textbf{Class template logical\_\allowbreak{}and.} Specifies the facility family Class template logical\_\allowbreak{}and. Read its declarations together with mandates, preconditions, effects, results, exceptions, complexity, and lifetime rules.
\item[22.10.10.3] \textbf{Class template logical\_\allowbreak{}or.} Specifies the facility family Class template logical\_\allowbreak{}or. Read its declarations together with mandates, preconditions, effects, results, exceptions, complexity, and lifetime rules.
\item[22.10.10.4] \textbf{Class template logical\_\allowbreak{}not.} Specifies the facility family Class template logical\_\allowbreak{}not. Read its declarations together with mandates, preconditions, effects, results, exceptions, complexity, and lifetime rules.
\item[22.10.11.1] \textbf{General.} Frames the terminology, applicability, and relationships needed to interpret General; detailed rules remain in the following subclauses.
\item[22.10.11.2] \textbf{Class template bit\_\allowbreak{}and.} Specifies the facility family Class template bit\_\allowbreak{}and. Read its declarations together with mandates, preconditions, effects, results, exceptions, complexity, and lifetime rules.
\item[22.10.11.3] \textbf{Class template bit\_\allowbreak{}or.} Specifies the facility family Class template bit\_\allowbreak{}or. Read its declarations together with mandates, preconditions, effects, results, exceptions, complexity, and lifetime rules.
\item[22.10.11.4] \textbf{Class template bit\_\allowbreak{}xor.} Specifies the facility family Class template bit\_\allowbreak{}xor. Read its declarations together with mandates, preconditions, effects, results, exceptions, complexity, and lifetime rules.
\item[22.10.11.5] \textbf{Class template bit\_\allowbreak{}not.} Specifies the facility family Class template bit\_\allowbreak{}not. Read its declarations together with mandates, preconditions, effects, results, exceptions, complexity, and lifetime rules.
\item[22.10.15.1] \textbf{General.} Frames the terminology, applicability, and relationships needed to interpret General; detailed rules remain in the following subclauses.
\item[22.10.15.2] \textbf{Class template is\_\allowbreak{}bind\_\allowbreak{}expression.} Specifies the facility family Class template is\_\allowbreak{}bind\_\allowbreak{}expression. Read its declarations together with mandates, preconditions, effects, results, exceptions, complexity, and lifetime rules.
\item[22.10.15.3] \textbf{Class template is\_\allowbreak{}placeholder.} Specifies the facility family Class template is\_\allowbreak{}placeholder. Read its declarations together with mandates, preconditions, effects, results, exceptions, complexity, and lifetime rules.
\item[22.10.15.4] \textbf{Function template bind.} Specifies the facility family Function template bind. Read its declarations together with mandates, preconditions, effects, results, exceptions, complexity, and lifetime rules.
\item[22.10.15.5] \textbf{Placeholders.} Specifies the facility family Placeholders. Read its declarations together with mandates, preconditions, effects, results, exceptions, complexity, and lifetime rules.
\item[22.10.17.1] \textbf{General.} Frames the terminology, applicability, and relationships needed to interpret General; detailed rules remain in the following subclauses.
\item[22.10.17.2] \textbf{Class bad\_\allowbreak{}function\_\allowbreak{}call.} Specifies the facility family Class bad\_\allowbreak{}function\_\allowbreak{}call. Read its declarations together with mandates, preconditions, effects, results, exceptions, complexity, and lifetime rules.
\item[22.10.17.3] \textbf{Class template function.} Specifies the facility family Class template function. Read its declarations together with mandates, preconditions, effects, results, exceptions, complexity, and lifetime rules.
\item[22.10.17.3.1] \textbf{General.} Frames the terminology, applicability, and relationships needed to interpret General; detailed rules remain in the following subclauses.
\item[22.10.17.3.2] \textbf{Constructors and destructor.} Specifies how Constructors and destructor establishes object state, including participation constraints, initialization effects, and exceptional exits.
\item[22.10.17.3.3] \textbf{Modifiers.} Specifies the facility family Modifiers. Read its declarations together with mandates, preconditions, effects, results, exceptions, complexity, and lifetime rules.
\item[22.10.17.3.4] \textbf{Capacity.} Specifies the facility family Capacity. Read its declarations together with mandates, preconditions, effects, results, exceptions, complexity, and lifetime rules.
\item[22.10.17.3.5] \textbf{Invocation.} Specifies the facility family Invocation. Read its declarations together with mandates, preconditions, effects, results, exceptions, complexity, and lifetime rules.
\item[22.10.17.3.6] \textbf{Target access.} Specifies the facility family Target access. Read its declarations together with mandates, preconditions, effects, results, exceptions, complexity, and lifetime rules.
\item[22.10.17.3.7] \textbf{Null pointer comparison operator functions.} Specifies the facility family Null pointer comparison operator functions. Read its declarations together with mandates, preconditions, effects, results, exceptions, complexity, and lifetime rules.
\item[22.10.17.3.8] \textbf{Specialized algorithms.} Defines the observable result of Specialized algorithms together with applicable iterator-domain, ordering, complexity, invalidation, and exception conditions.
\item[22.10.17.4] \textbf{Move only wrapper.} Specifies the facility family Move only wrapper. Read its declarations together with mandates, preconditions, effects, results, exceptions, complexity, and lifetime rules.
\item[22.10.17.4.1] \textbf{General.} Frames the terminology, applicability, and relationships needed to interpret General; detailed rules remain in the following subclauses.
\item[22.10.17.4.2] \textbf{Class template move\_\allowbreak{}only\_\allowbreak{}function.} Specifies the facility family Class template move\_\allowbreak{}only\_\allowbreak{}function. Read its declarations together with mandates, preconditions, effects, results, exceptions, complexity, and lifetime rules.
\item[22.10.17.4.3] \textbf{Constructors, assignment, and destructor.} Specifies how Constructors, assignment, and destructor establishes object state, including participation constraints, initialization effects, and exceptional exits.
\item[22.10.17.4.4] \textbf{Invocation.} Specifies the facility family Invocation. Read its declarations together with mandates, preconditions, effects, results, exceptions, complexity, and lifetime rules.
\item[22.10.17.4.5] \textbf{Utility.} Specifies the facility family Utility. Read its declarations together with mandates, preconditions, effects, results, exceptions, complexity, and lifetime rules.
\item[22.10.18.1] \textbf{General.} Frames the terminology, applicability, and relationships needed to interpret General; detailed rules remain in the following subclauses.
\item[22.10.18.2] \textbf{Class template default\_\allowbreak{}searcher.} Specifies the facility family Class template default\_\allowbreak{}searcher. Read its declarations together with mandates, preconditions, effects, results, exceptions, complexity, and lifetime rules.
\item[22.10.18.3] \textbf{Class template boyer\_\allowbreak{}moore\_\allowbreak{}searcher.} Specifies the facility family Class template boyer\_\allowbreak{}moore\_\allowbreak{}searcher. Read its declarations together with mandates, preconditions, effects, results, exceptions, complexity, and lifetime rules.
\item[22.10.18.4] \textbf{Class template boyer\_\allowbreak{}moore\_\allowbreak{}horspool\_\allowbreak{}searcher.} Specifies the facility family Class template boyer\_\allowbreak{}moore\_\allowbreak{}horspool\_\allowbreak{}searcher. Read its declarations together with mandates, preconditions, effects, results, exceptions, complexity, and lifetime rules.
\end{description}
\Needspace{8\baselineskip}
\subsection*{22.11 Class type\_\allowbreak{}index}
\begin{longtable}{@{}>{\RaggedRight\arraybackslash\columncolor{CornellReferenceCueGray}}p{0.28\textwidth}>{\RaggedRight\arraybackslash}p{0.64\textwidth}@{}}
\rowcolor{CornellReferenceNavy}\color{white}\textbf{Cue / recall prompt} & \color{white}\textbf{Notes}\\\endfirsthead
\rowcolor{CornellReferenceNavy}\color{white}\textbf{Cue / recall prompt} & \color{white}\textbf{Notes (continued)}\\\endhead
\arrayrulecolor{CornellReferenceLineGray}
\textbf{What contract governs Header <typeindex> synopsis?}\\[-1mm]{\scriptsize\CornellClauseRef{22.11.1}} & Catalogs the declarations made available by Header <typeindex> synopsis. The synopsis is an interface map; the detailed specifications supply constraints and behavior.\\\hline
\textbf{What contract governs type\_\allowbreak{}index overview?}\\[-1mm]{\scriptsize\CornellClauseRef{22.11.2}} & Explains the traversal or view contract for type\_\allowbreak{}index overview; prove domain validity, lifetime, sentinel compatibility, and any borrowing or invalidation property.\\\hline
\textbf{What contract governs type\_\allowbreak{}index members?}\\[-1mm]{\scriptsize\CornellClauseRef{22.11.3}} & Specifies the facility family type\_\allowbreak{}index members. Read its declarations together with mandates, preconditions, effects, results, exceptions, complexity, and lifetime rules.\\\hline
\textbf{What contract governs Hash support?}\\[-1mm]{\scriptsize\CornellClauseRef{22.11.4}} & Specifies the facility family Hash support. Read its declarations together with mandates, preconditions, effects, results, exceptions, complexity, and lifetime rules.\\\hline
\end{longtable}
\begin{CornellReferenceSummaryBox}Class type\_\allowbreak{}index supplies the library semantics portion of this clause. The key review move is to connect its local wording to the clause-wide model, then classify each condition by normative force before relying on it.\end{CornellReferenceSummaryBox}
\Needspace{8\baselineskip}
\subsection*{22.12 Execution policies}
\begin{longtable}{@{}>{\RaggedRight\arraybackslash\columncolor{CornellReferenceCueGray}}p{0.28\textwidth}>{\RaggedRight\arraybackslash}p{0.64\textwidth}@{}}
\rowcolor{CornellReferenceNavy}\color{white}\textbf{Cue / recall prompt} & \color{white}\textbf{Notes}\\\endfirsthead
\rowcolor{CornellReferenceNavy}\color{white}\textbf{Cue / recall prompt} & \color{white}\textbf{Notes (continued)}\\\endhead
\arrayrulecolor{CornellReferenceLineGray}
\textbf{What contract governs General?}\\[-1mm]{\scriptsize\CornellClauseRef{22.12.1}} & Frames the terminology, applicability, and relationships needed to interpret General; detailed rules remain in the following subclauses.\\\hline
\textbf{What contract governs Header <execution> synopsis?}\\[-1mm]{\scriptsize\CornellClauseRef{22.12.2}} & Catalogs the declarations made available by Header <execution> synopsis. The synopsis is an interface map; the detailed specifications supply constraints and behavior.\\\hline
\textbf{What contract governs Execution policy type trait?}\\[-1mm]{\scriptsize\CornellClauseRef{22.12.3}} & Specifies the facility family Execution policy type trait. Read its declarations together with mandates, preconditions, effects, results, exceptions, complexity, and lifetime rules.\\\hline
\textbf{What contract governs Sequenced execution policy?}\\[-1mm]{\scriptsize\CornellClauseRef{22.12.4}} & Specifies the facility family Sequenced execution policy. Read its declarations together with mandates, preconditions, effects, results, exceptions, complexity, and lifetime rules.\\\hline
\textbf{What contract governs Parallel execution policy?}\\[-1mm]{\scriptsize\CornellClauseRef{22.12.5}} & Specifies the facility family Parallel execution policy. Read its declarations together with mandates, preconditions, effects, results, exceptions, complexity, and lifetime rules.\\\hline
\textbf{What contract governs Parallel and unsequenced execution policy?}\\[-1mm]{\scriptsize\CornellClauseRef{22.12.6}} & Specifies the facility family Parallel and unsequenced execution policy. Read its declarations together with mandates, preconditions, effects, results, exceptions, complexity, and lifetime rules.\\\hline
\textbf{What contract governs Unsequenced execution policy?}\\[-1mm]{\scriptsize\CornellClauseRef{22.12.7}} & Specifies the facility family Unsequenced execution policy. Read its declarations together with mandates, preconditions, effects, results, exceptions, complexity, and lifetime rules.\\\hline
\textbf{What contract governs Execution policy objects?}\\[-1mm]{\scriptsize\CornellClauseRef{22.12.8}} & Specifies the facility family Execution policy objects. Read its declarations together with mandates, preconditions, effects, results, exceptions, complexity, and lifetime rules.\\\hline
\end{longtable}
\begin{CornellReferenceSummaryBox}Execution policies supplies the library semantics portion of this clause. The key review move is to connect its local wording to the clause-wide model, then classify each condition by normative force before relying on it.\end{CornellReferenceSummaryBox}
\Needspace{8\baselineskip}
\subsection*{22.13 Primitive numeric conversions}
\begin{longtable}{@{}>{\RaggedRight\arraybackslash\columncolor{CornellReferenceCueGray}}p{0.28\textwidth}>{\RaggedRight\arraybackslash}p{0.64\textwidth}@{}}
\rowcolor{CornellReferenceNavy}\color{white}\textbf{Cue / recall prompt} & \color{white}\textbf{Notes}\\\endfirsthead
\rowcolor{CornellReferenceNavy}\color{white}\textbf{Cue / recall prompt} & \color{white}\textbf{Notes (continued)}\\\endhead
\arrayrulecolor{CornellReferenceLineGray}
\textbf{What contract governs Header <charconv> synopsis?}\\[-1mm]{\scriptsize\CornellClauseRef{22.13.1}} & Catalogs the declarations made available by Header <charconv> synopsis. The synopsis is an interface map; the detailed specifications supply constraints and behavior.\\\hline
\textbf{When is Primitive numeric output conversion permitted?}\\[-1mm]{\scriptsize\CornellClauseRef{22.13.2}} & Specifies source and destination conditions for Primitive numeric output conversion, the resulting type or value category, and any information-loss, pointer, qualification, or lifetime consequence.\\\hline
\textbf{When is Primitive numeric input conversion permitted?}\\[-1mm]{\scriptsize\CornellClauseRef{22.13.3}} & Specifies source and destination conditions for Primitive numeric input conversion, the resulting type or value category, and any information-loss, pointer, qualification, or lifetime consequence.\\\hline
\end{longtable}
\begin{CornellReferenceSummaryBox}Primitive numeric conversions supplies the library semantics portion of this clause. The key review move is to connect its local wording to the clause-wide model, then classify each condition by normative force before relying on it.\end{CornellReferenceSummaryBox}
\Needspace{8\baselineskip}
\subsection*{22.14 Formatting}
\begin{longtable}{@{}>{\RaggedRight\arraybackslash\columncolor{CornellReferenceCueGray}}p{0.28\textwidth}>{\RaggedRight\arraybackslash}p{0.64\textwidth}@{}}
\rowcolor{CornellReferenceNavy}\color{white}\textbf{Cue / recall prompt} & \color{white}\textbf{Notes}\\\endfirsthead
\rowcolor{CornellReferenceNavy}\color{white}\textbf{Cue / recall prompt} & \color{white}\textbf{Notes (continued)}\\\endhead
\arrayrulecolor{CornellReferenceLineGray}
\textbf{What contract governs Header <format> synopsis?}\\[-1mm]{\scriptsize\CornellClauseRef{22.14.1}} & Catalogs the declarations made available by Header <format> synopsis. The synopsis is an interface map; the detailed specifications supply constraints and behavior.\\\hline
\textbf{What contract governs Format string?}\\[-1mm]{\scriptsize\CornellClauseRef{22.14.2}} & Specifies text-sequence behavior for Format string; establish character type, extent, ownership, null termination, encoding assumptions, and invalidation before use.\\\hline
\textbf{What contract governs Error reporting?}\\[-1mm]{\scriptsize\CornellClauseRef{22.14.3}} & Defines the failure model for Error reporting, including how an error is represented, reported, propagated, or converted into a diagnostic consequence.\\\hline
\textbf{What contract governs Class template basic\_\allowbreak{}format\_\allowbreak{}string?}\\[-1mm]{\scriptsize\CornellClauseRef{22.14.4}} & Specifies text-sequence behavior for Class template basic\_\allowbreak{}format\_\allowbreak{}string; establish character type, extent, ownership, null termination, encoding assumptions, and invalidation before use.\\\hline
\textbf{What contract governs Formatting functions?}\\[-1mm]{\scriptsize\CornellClauseRef{22.14.5}} & Specifies the facility family Formatting functions. Read its declarations together with mandates, preconditions, effects, results, exceptions, complexity, and lifetime rules.\\\hline
\textbf{What contract governs Formatter?}\\[-1mm]{\scriptsize\CornellClauseRef{22.14.6}} & Specifies the facility family Formatter. Read its declarations together with mandates, preconditions, effects, results, exceptions, complexity, and lifetime rules.\\\hline
\textbf{What contract governs Formatting of ranges?}\\[-1mm]{\scriptsize\CornellClauseRef{22.14.7}} & Explains the traversal or view contract for Formatting of ranges; prove domain validity, lifetime, sentinel compatibility, and any borrowing or invalidation property.\\\hline
\textbf{What contract governs Arguments?}\\[-1mm]{\scriptsize\CornellClauseRef{22.14.8}} & Specifies the facility family Arguments. Read its declarations together with mandates, preconditions, effects, results, exceptions, complexity, and lifetime rules.\\\hline
\textbf{What contract governs Tuple formatter?}\\[-1mm]{\scriptsize\CornellClauseRef{22.14.9}} & Specifies the facility family Tuple formatter. Read its declarations together with mandates, preconditions, effects, results, exceptions, complexity, and lifetime rules.\\\hline
\textbf{What contract governs Class format\_\allowbreak{}error?}\\[-1mm]{\scriptsize\CornellClauseRef{22.14.10}} & Defines the failure model for Class format\_\allowbreak{}error, including how an error is represented, reported, propagated, or converted into a diagnostic consequence.\\\hline
\end{longtable}
\begin{CornellReferenceSummaryBox}Formatting supplies the library semantics portion of this clause. The key review move is to connect its local wording to the clause-wide model, then classify each condition by normative force before relying on it.\end{CornellReferenceSummaryBox}
\Needspace{5\baselineskip}\paragraph{Detailed coverage ledger.}
\begin{description}[style=nextline,leftmargin=2.8cm,font=\normalfont\bfseries\color{CornellReferenceTeal}]
\item[22.14.2.1] \textbf{General.} Frames the terminology, applicability, and relationships needed to interpret General; detailed rules remain in the following subclauses.
\item[22.14.2.2] \textbf{Standard format specifiers.} Specifies the facility family Standard format specifiers. Read its declarations together with mandates, preconditions, effects, results, exceptions, complexity, and lifetime rules.
\item[22.14.6.1] \textbf{Formatter requirements.} States the admissibility and semantic obligations for Formatter requirements. Separate compile-time participation rules from caller preconditions and runtime effects.
\item[22.14.6.2] \textbf{Concept formattable.} Defines or applies the generic requirement Concept formattable; syntactic satisfaction must be paired with every stated semantic and equality-preservation expectation.
\item[22.14.6.3] \textbf{Formatter specializations.} Specifies the facility family Formatter specializations. Read its declarations together with mandates, preconditions, effects, results, exceptions, complexity, and lifetime rules.
\item[22.14.6.4] \textbf{Formatting escaped characters and strings.} Specifies text-sequence behavior for Formatting escaped characters and strings; establish character type, extent, ownership, null termination, encoding assumptions, and invalidation before use.
\item[22.14.6.5] \textbf{Class template basic\_\allowbreak{}format\_\allowbreak{}parse\_\allowbreak{}context.} Specifies the facility family Class template basic\_\allowbreak{}format\_\allowbreak{}parse\_\allowbreak{}context. Read its declarations together with mandates, preconditions, effects, results, exceptions, complexity, and lifetime rules.
\item[22.14.6.6] \textbf{Class template basic\_\allowbreak{}format\_\allowbreak{}context.} Specifies the facility family Class template basic\_\allowbreak{}format\_\allowbreak{}context. Read its declarations together with mandates, preconditions, effects, results, exceptions, complexity, and lifetime rules.
\item[22.14.7.1] \textbf{Variable template format\_\allowbreak{}kind.} Specifies the facility family Variable template format\_\allowbreak{}kind. Read its declarations together with mandates, preconditions, effects, results, exceptions, complexity, and lifetime rules.
\item[22.14.7.2] \textbf{Class template range\_\allowbreak{}formatter.} Explains the traversal or view contract for Class template range\_\allowbreak{}formatter; prove domain validity, lifetime, sentinel compatibility, and any borrowing or invalidation property.
\item[22.14.7.3] \textbf{Class template range-default-formatter.} Explains the traversal or view contract for Class template range-default-formatter; prove domain validity, lifetime, sentinel compatibility, and any borrowing or invalidation property.
\item[22.14.7.4] \textbf{Specialization of range-default-formatter for maps.} Explains the traversal or view contract for Specialization of range-default-formatter for maps; prove domain validity, lifetime, sentinel compatibility, and any borrowing or invalidation property.
\item[22.14.7.5] \textbf{Specialization of range-default-formatter for sets.} Explains the traversal or view contract for Specialization of range-default-formatter for sets; prove domain validity, lifetime, sentinel compatibility, and any borrowing or invalidation property.
\item[22.14.7.6] \textbf{Specialization of range-default-formatter for strings.} Explains the traversal or view contract for Specialization of range-default-formatter for strings; prove domain validity, lifetime, sentinel compatibility, and any borrowing or invalidation property.
\item[22.14.8.1] \textbf{Class template basic\_\allowbreak{}format\_\allowbreak{}arg.} Specifies the facility family Class template basic\_\allowbreak{}format\_\allowbreak{}arg. Read its declarations together with mandates, preconditions, effects, results, exceptions, complexity, and lifetime rules.
\item[22.14.8.2] \textbf{Class template format-arg-store.} Specifies the facility family Class template format-arg-store. Read its declarations together with mandates, preconditions, effects, results, exceptions, complexity, and lifetime rules.
\item[22.14.8.3] \textbf{Class template basic\_\allowbreak{}format\_\allowbreak{}args.} Specifies the facility family Class template basic\_\allowbreak{}format\_\allowbreak{}args. Read its declarations together with mandates, preconditions, effects, results, exceptions, complexity, and lifetime rules.
\end{description}
\Needspace{8\baselineskip}
\subsection*{22.15 Bit manipulation}
\begin{longtable}{@{}>{\RaggedRight\arraybackslash\columncolor{CornellReferenceCueGray}}p{0.28\textwidth}>{\RaggedRight\arraybackslash}p{0.64\textwidth}@{}}
\rowcolor{CornellReferenceNavy}\color{white}\textbf{Cue / recall prompt} & \color{white}\textbf{Notes}\\\endfirsthead
\rowcolor{CornellReferenceNavy}\color{white}\textbf{Cue / recall prompt} & \color{white}\textbf{Notes (continued)}\\\endhead
\arrayrulecolor{CornellReferenceLineGray}
\textbf{What contract governs General?}\\[-1mm]{\scriptsize\CornellClauseRef{22.15.1}} & Frames the terminology, applicability, and relationships needed to interpret General; detailed rules remain in the following subclauses.\\\hline
\textbf{What contract governs Header <bit> synopsis?}\\[-1mm]{\scriptsize\CornellClauseRef{22.15.2}} & Catalogs the declarations made available by Header <bit> synopsis. The synopsis is an interface map; the detailed specifications supply constraints and behavior.\\\hline
\textbf{What contract governs Function template bit\_\allowbreak{}cast?}\\[-1mm]{\scriptsize\CornellClauseRef{22.15.3}} & Specifies the facility family Function template bit\_\allowbreak{}cast. Read its declarations together with mandates, preconditions, effects, results, exceptions, complexity, and lifetime rules.\\\hline
\textbf{What contract governs byteswap?}\\[-1mm]{\scriptsize\CornellClauseRef{22.15.4}} & Governs state transfer or exchange for byteswap; check self-operation, validity after movement, exception behavior, and invalidation.\\\hline
\textbf{What contract governs Integral powers of 2?}\\[-1mm]{\scriptsize\CornellClauseRef{22.15.5}} & Specifies the facility family Integral powers of 2. Read its declarations together with mandates, preconditions, effects, results, exceptions, complexity, and lifetime rules.\\\hline
\textbf{What contract governs Rotating?}\\[-1mm]{\scriptsize\CornellClauseRef{22.15.6}} & Specifies the facility family Rotating. Read its declarations together with mandates, preconditions, effects, results, exceptions, complexity, and lifetime rules.\\\hline
\textbf{What contract governs Counting?}\\[-1mm]{\scriptsize\CornellClauseRef{22.15.7}} & Specifies the facility family Counting. Read its declarations together with mandates, preconditions, effects, results, exceptions, complexity, and lifetime rules.\\\hline
\textbf{What contract governs Endian?}\\[-1mm]{\scriptsize\CornellClauseRef{22.15.8}} & Specifies the facility family Endian. Read its declarations together with mandates, preconditions, effects, results, exceptions, complexity, and lifetime rules.\\\hline
\end{longtable}
\begin{CornellReferenceSummaryBox}Bit manipulation supplies the library semantics portion of this clause. The key review move is to connect its local wording to the clause-wide model, then classify each condition by normative force before relying on it.\end{CornellReferenceSummaryBox}
\section{Language-Lawyer and Library-Usage Pitfalls}
\begin{CornellReferencePitfallBox}\begin{itemize}[label=\textcolor{CornellReferenceAmber}{$\blacktriangleright$}]
\item Reading a synopsis as the complete behavioral contract.
\item Ignoring mandates, preconditions, exception behavior, or complexity requirements.
\item Using an exposition-only name as though it were public API.
\item Assuming invalidation, ownership, or thread-safety properties that are not stated.
\item Treating successful compilation as proof that semantic requirements and preconditions are met.
\end{itemize}\end{CornellReferencePitfallBox}
\section{Programmer and Implementation Checklist}
\begin{CornellReferenceChecklistBox}\begin{itemize}[label=$\square$]
\item Identify the declaring header or module exposure for each facility.
\item Check template arguments and mandates before evaluating an operation.
\item Establish every precondition at the call site.
\item Record effects, returned value, postconditions, and exceptions separately.
\item Audit iterator, reference, pointer, and object lifetime after mutation.
\item Verify stated complexity and synchronization assumptions.
\item For \CornellClauseRef{22.1}, verify general against the precise rule category and all cited dependencies.
\item For \CornellClauseRef{22.2}, verify utility components against the precise rule category and all cited dependencies.
\item For \CornellClauseRef{22.3}, verify pairs against the precise rule category and all cited dependencies.
\item For \CornellClauseRef{22.4}, verify tuples against the precise rule category and all cited dependencies.
\end{itemize}\end{CornellReferenceChecklistBox}
\section{Clause Synthesis}
\begin{CornellReferenceOrientationBox}Collects broadly reusable facilities: utilities, tuples, sum types, expected values, bitsets, function wrappers, binders, hashing, type indexing, clocks and execution support, formatting, and other general components. A cross-cutting toolbox shared by applications and other library clauses. A sound review distinguishes syntax and participation from runtime preconditions, keeps notes and examples non-mandatory, and follows cross-clause dependencies before making a portability claim.\end{CornellReferenceOrientationBox}
\section{Self-Test Questions}
\begin{enumerate}
\item What role does General play in the clause's overall model?
\item Which normative distinctions must be kept separate when applying Utility components?
\item How would you audit a program or implementation that depends on Pairs?
\item Compare Tuples with Optional objects; what different question does each answer?
\item What role does Optional objects play in the clause's overall model?
\item Which normative distinctions must be kept separate when applying Variants?
\item How would you audit a program or implementation that depends on Storage for any type?
\item Compare Expected objects with Bitsets; what different question does each answer?
\item What role does Bitsets play in the clause's overall model?
\item Which normative distinctions must be kept separate when applying Function objects?
\item How would you audit a program or implementation that depends on Class type\_\allowbreak{}index?
\item Compare Execution policies with Primitive numeric conversions; what different question does each answer?
\item What role does Primitive numeric conversions play in the clause's overall model?
\item Which normative distinctions must be kept separate when applying Formatting?
\item Scenario: a program relies on an unstated implementation behavior while using Bit manipulation. What must be checked before calling the program portable?
\item Scenario: a program relies on an unstated implementation behavior while using Header <utility> synopsis. What must be checked before calling the program portable?
\end{enumerate}
\clearpage\section{Answer Key}
\begin{enumerate}
\item General is one part of the clause's library semantics model. Its role is understood by connecting the local rule in 22.15.1 to the clause orientation and its dependent concepts.
\item Keep formation and well-formedness separate from runtime preconditions and effects. Also distinguish mandatory wording from notes, implementation choice from undefined behavior, and interface declarations from detailed contracts; see 22.2.
\item Audit Pairs by locating the controlling wording in 22.3, listing inputs and type constraints, proving any preconditions, recording effects and failure behavior, and checking lifetime, invalidation, complexity, and synchronization where applicable.
\item Tuples and Optional objects occupy different steps or facility groups. Analyze each under its own subclause, then follow the explicit dependency between them rather than transferring one set of guarantees to the other; begin at 22.4.
\item Optional objects is one part of the clause's library semantics model. Its role is understood by connecting the local rule in 22.5 to the clause orientation and its dependent concepts.
\item Keep formation and well-formedness separate from runtime preconditions and effects. Also distinguish mandatory wording from notes, implementation choice from undefined behavior, and interface declarations from detailed contracts; see 22.6.
\item Audit Storage for any type by locating the controlling wording in 22.7, listing inputs and type constraints, proving any preconditions, recording effects and failure behavior, and checking lifetime, invalidation, complexity, and synchronization where applicable.
\item Expected objects and Bitsets occupy different steps or facility groups. Analyze each under its own subclause, then follow the explicit dependency between them rather than transferring one set of guarantees to the other; begin at 22.8.
\item Bitsets is one part of the clause's library semantics model. Its role is understood by connecting the local rule in 22.9 to the clause orientation and its dependent concepts.
\item Keep formation and well-formedness separate from runtime preconditions and effects. Also distinguish mandatory wording from notes, implementation choice from undefined behavior, and interface declarations from detailed contracts; see 22.10.
\item Audit Class type\_\allowbreak{}index by locating the controlling wording in 22.11, listing inputs and type constraints, proving any preconditions, recording effects and failure behavior, and checking lifetime, invalidation, complexity, and synchronization where applicable.
\item Execution policies and Primitive numeric conversions occupy different steps or facility groups. Analyze each under its own subclause, then follow the explicit dependency between them rather than transferring one set of guarantees to the other; begin at 22.12.
\item Primitive numeric conversions is one part of the clause's library semantics model. Its role is understood by connecting the local rule in 22.13 to the clause orientation and its dependent concepts.
\item Keep formation and well-formedness separate from runtime preconditions and effects. Also distinguish mandatory wording from notes, implementation choice from undefined behavior, and interface declarations from detailed contracts; see 22.14.
\item First identify the exact requirement in 22.15, then classify the relied-on behavior as required, unspecified, implementation-defined, conditionally supported, or undefined. Portability requires satisfying the program obligations and avoiding dependence on a choice not guaranteed in the target environments.
\item First identify the exact requirement in 22.2.1, then classify the relied-on behavior as required, unspecified, implementation-defined, conditionally supported, or undefined. Portability requires satisfying the program obligations and avoiding dependence on a choice not guaranteed in the target environments.
\end{enumerate}
\section{Key-Term Recap}
\begin{longtable}{@{}>{\RaggedRight\arraybackslash}p{0.28\textwidth}>{\RaggedRight\arraybackslash}p{0.64\textwidth}@{}}\toprule\textbf{Term} & \textbf{Working meaning}\\\midrule\endfirsthead\toprule\textbf{Term} & \textbf{Working meaning (continued)}\\\midrule\endhead\bottomrule\endfoot
pair & A product type with two elements.\\
tuple & A fixed-size heterogeneous product type.\\
optional & A value type representing zero or one contained objects.\\
variant & A type-safe discriminated union.\\
expected & A value-or-error result type.\\
function & A polymorphic callable wrapper.\\
hash & A customization family producing hash values.\\
format & A type-safe formatting facility governed by format strings and formatter specializations.\\
\end{longtable}
\CornellTakeaway{General utilities reduce boilerplate by giving common value, callable, sum-type, and formatting patterns precise contracts.}
\end{document}
